% StormerRule.tex
% Mathematical derivation of Störmer's modified midpoint rule for
% second-order ODEs of the form y'' = f(x, y).
% Reference: NR3 §17.4; Gragg (1965); Hairer, Nørsett, Wanner Vol. I §§II.14, III.10.

\documentclass[11pt]{article}

\usepackage{amsmath, amssymb, amsthm}
\usepackage{mathtools}
\usepackage{geometry}
\usepackage{hyperref}
\usepackage{booktabs}
\usepackage{enumitem}

\geometry{margin=1.2in}

\newtheorem{theorem}{Theorem}
\newtheorem{corollary}{Corollary}
\newtheorem{remark}{Remark}

\DeclareMathOperator{\OO}{\mathcal{O}}

\title{\textbf{Störmer's Rule}\\
       \large Mathematical Derivation from First Principles}
\author{Cosmos GNC Library --- Propulsion Repository}
\date{2026}

\begin{document}
\maketitle
\tableofcontents
\bigskip

% ---------------------------------------------------------------------------
\section{Problem Statement}
% ---------------------------------------------------------------------------

We consider a second-order ordinary differential equation of the form
\begin{equation}
  \label{eq:stormer-ode}
  y'' = f(x,\, y),
  \qquad y(x_0) = y_0,\quad y'(x_0) = v_0,
\end{equation}
where $y : \mathbb{R} \to \mathbb{R}^n$, $f : \mathbb{R} \times \mathbb{R}^n \to
\mathbb{R}^n$, and crucially \emph{$f$ does not depend on $y'$}.

\begin{remark}[Physical motivation]
  Equation~\eqref{eq:stormer-ode} describes any purely position-dependent force:
  gravitational $n$-body, Keplerian orbits, conservative mechanical systems, and
  the Störmer problem (charged particle in a magnetic field) that originally
  motivated the method.  For the two-body orbital problem:
  $f(\mathbf{r}) = -\mu\,\mathbf{r}/|\mathbf{r}|^3$.
\end{remark}

\paragraph{Goal.}
Advance the solution from $x_0$ to $x_0 + H$ by subdividing $H$ into $M$ equal
sub-steps of size $h = H/M$, in a way that:
\begin{enumerate}
  \item exploits the absence of $y'$ in $f$,
  \item produces an \emph{even} error series in $h$ (enabling Bulirsch-Stoer
        Richardson extrapolation), and
  \item requires only $M/2$ new evaluations of $f$ per column (half the cost of
        applying the standard modified midpoint rule to the equivalent first-order
        system).
\end{enumerate}

% ---------------------------------------------------------------------------
\section{The Störmer Three-Point Recurrence}
\label{sec:stormer-3pt}
% ---------------------------------------------------------------------------

\subsection{Derivation from Taylor's Theorem}

Write the Taylor expansions of $y$ centered at $x_m = x_0 + mh$:
\begin{align}
  y(x_m + h) &= y_m + h y'_m + \tfrac{h^2}{2} y''_m
                + \tfrac{h^3}{6} y'''_m
                + \tfrac{h^4}{24} y^{(4)}_m + \OO(h^5), \label{eq:taylor-fwd}\\
  y(x_m - h) &= y_m - h y'_m + \tfrac{h^2}{2} y''_m
                - \tfrac{h^3}{6} y'''_m
                + \tfrac{h^4}{24} y^{(4)}_m + \OO(h^5). \label{eq:taylor-bwd}
\end{align}
Adding \eqref{eq:taylor-fwd} and \eqref{eq:taylor-bwd} cancels all odd powers:
\begin{equation}
  y(x_m + h) + y(x_m - h) = 2y_m + h^2 y''_m + \tfrac{h^4}{12} y^{(4)}_m + \OO(h^6).
\end{equation}
Substituting $y''_m = f(x_m, y_m)$ and rearranging:
\begin{equation}
  \label{eq:stormer-3pt}
  \boxed{y_{m+1} = 2y_m - y_{m-1} + h^2 f(x_m,\, y_m) + \OO(h^4).}
\end{equation}
This is the \textbf{Störmer three-point formula}: given two successive solution
values $y_{m-1}$ and $y_m$, the next value $y_{m+1}$ requires a single evaluation
of $f$ and is accurate to $\OO(h^4)$ locally.

\begin{remark}[No velocity in the recurrence]
  Equation~\eqref{eq:stormer-3pt} uses $y_{m-1}$, $y_m$, and $f(x_m, y_m)$ only.
  The velocity $y'$ has been eliminated algebraically.  The price is the need for
  two starting values.
\end{remark}

\subsection{The \texorpdfstring{$\delta$}{delta}-Increment Formulation}

Define the position increment $\delta_m \coloneqq \theta_{m+1} - \theta_m$, the change in the numerical solution over one sub-step.  From
\eqref{eq:stormer-3pt}:
\begin{align}
  \delta_m &= y_{m+1} - y_m \notag \\
           &= (2y_m - y_{m-1} + h^2 f_m) - y_m \notag \\
           &= (y_m - y_{m-1}) + h^2 f_m \notag \\
           &= \delta_{m-1} + h^2 f_m. \label{eq:delta-recurrence}
\end{align}
This telescoping relation propagates $\delta$ forward with only $h^2 f_m$ added at
each step.  The velocity information is encoded in the initial $\delta_0$.

\paragraph{Starting value from Taylor:}
The value $y_1 = y(x_0 + h)$ from \eqref{eq:taylor-fwd} with $m=0$:
\begin{equation}
  \label{eq:y1}
  y_1 = y_0 + h v_0 + \tfrac{h^2}{2} f_0 + \OO(h^3),
\end{equation}
so
\begin{equation}
  \label{eq:delta0}
  \boxed{\delta_0 = h\, v_0 + \tfrac{h^2}{2}\, f_0.}
\end{equation}

% ---------------------------------------------------------------------------
\section{Order Analysis: What ``Order $p$'' Means}
\label{sec:order}
% ---------------------------------------------------------------------------

The three-point recurrence~\eqref{eq:stormer-3pt} was derived from Taylor's
theorem.  We now ask precisely: \emph{how fast does the global error shrink as
$h \to 0$?}  The answer, ``order $p$'', refers entirely to powers of $h$ --- it
is \textbf{not} the degree of a polynomial fitted to $f$.

\subsection{The HNW Framework}

Following Hairer, N{\o}rsett \& Wanner~\cite{HNW} (henceforth \textbf{HNW}),
write the general second-order multistep method for $y'' = f(x,y)$ as
\begin{equation}
  \label{eq:general-multistep}
  \sum_{i=0}^{k} \alpha_i\, y_{n+i}
  \;=\; h^2 \sum_{i=0}^{k} \beta_i\, f(x_{n+i},\, y_{n+i}),
  \qquad \text{HNW (III.10, Eq.~10.19)}
\end{equation}
with scalar sequences $\{\alpha_i\}$ and $\{\beta_i\}$, $\alpha_k \neq 0$.
Define the two \textbf{characteristic polynomials}
\begin{equation}
  \label{eq:char-polys}
  \varrho(\zeta) = \sum_{i=0}^{k} \alpha_i\,\zeta^i,
  \qquad
  \sigma(\zeta) = \sum_{i=0}^{k} \beta_i\,\zeta^i.
\end{equation}
These encode the entire method: $\varrho$ governs stability and $\sigma$ governs
accuracy.

The residual when the \emph{exact} solution is substituted into the method is
measured by the \textbf{linear difference operator}
(HNW, Eq.~10.23, p.~467):
\begin{equation}
  \label{eq:L-operator}
  \mathcal{L}[y,\, x,\, h]
  \;=\; \varrho(E)\,y(x) \;-\; h^2\,\sigma(E)\,y''(x)
  \;=\; \sum_{i=0}^{k}
        \Bigl[\alpha_i\, y(x + ih) - h^2\,\beta_i\, y''(x + ih)\Bigr],
\end{equation}
where $E$ is the shift operator $E\,y(x) = y(x+h)$.  Along the exact solution
$y''(x+ih) = f\bigl(x+ih,\, y(x+ih)\bigr)$, so $\mathcal{L}$ is a pure
function of the smooth exact solution $y$ and the step $h$.

\begin{theorem}[HNW Definition~10.2 and Theorem~10.3, p.~467--468]
\label{thm:order}
Method~\eqref{eq:general-multistep} is \textbf{consistent of order $p$} if
\begin{equation}
  \label{eq:order-condition}
  \mathcal{L}[y,\, x,\, h] = \OO(h^{p+2})
  \quad \text{for every sufficiently smooth } y.
\end{equation}
Three equivalent conditions (HNW Theorem~10.3, p.~468) are:
\begin{enumerate}[label=(\roman*)]
  \item \emph{Moment form.}
    $\displaystyle\sum_i \alpha_i = 0$,\;
    $\displaystyle\sum_i i\,\alpha_i = 0$,\; and for $q = 2, \ldots, p{+}1$:
    \begin{equation}
      \label{eq:moment-conditions}
      \sum_{i=0}^{k} \alpha_i\, i^q
      \;=\; q(q-1)\sum_{i=0}^{k} \beta_i\, i^{q-2}.
    \end{equation}
  \item \emph{Generating-polynomial form.}
    \begin{equation}
      \label{eq:gen-poly-order}
      \varrho(e^h) - h^2\,\sigma(e^h) = \OO(h^{p+2})
      \quad\text{as } h \to 0.
    \end{equation}
  \item \emph{Root form near $\zeta = 1$.}
    \begin{equation}
      \frac{\varrho(\zeta)}{(\log\zeta)^2} - \sigma(\zeta)
      = \OO\!\bigl((\zeta-1)^p\bigr)
      \quad\text{as } \zeta \to 1.
    \end{equation}
\end{enumerate}
\end{theorem}

\paragraph{Why $h^{p+2}$, not $h^p$?}
The method equates the $h^2$-scaled forcing $h^2 f$ to the finite-difference
$\sum \alpha_i y_{n+i}$, so the entire equation is naturally $\OO(h^2)$.  The
residual $\mathcal{L}$ is therefore already premultiplied by $h^2$; order $p$
means the \emph{extra} smallness beyond $h^2$ is $h^p$ more --- i.e.,
$h^{p+2}$ total.  The global error, obtained by accumulating $\OO(h^{p+2})$
local residuals over $\OO(1/h)$ steps and dividing once by the implicit $h^2$
factor, is $\OO(h^p)$ --- matching Definition~10.2.  This is confirmed by:

\begin{theorem}[HNW Theorem~10.6, p.~468 --- Global Convergence]
\label{thm:convergence}
If method~\eqref{eq:general-multistep} is stable and of order $p$, and the
starting values satisfy $y(x_j) - y_j = \OO(h^{p+1})$ for $j = 0,\ldots,k-1$,
then
\begin{equation}
  \label{eq:global-error}
  \bigl|y(x_n) - y_n\bigr| \leq C\,h^p
  \qquad \text{for all } 0 \leq hn \leq \mathrm{const}.
\end{equation}
\end{theorem}

\subsection{Verification for the \texorpdfstring{$k=2$}{k=2} Störmer Method}
\label{ssec:order-k2}

For the three-point formula~\eqref{eq:stormer-3pt}:
\begin{equation}
  \label{eq:k2-coefficients}
  \alpha_0 = 1,\;\alpha_1 = -2,\;\alpha_2 = 1;
  \qquad
  \beta_0 = 0,\;\beta_1 = 1,\;\beta_2 = 0.
\end{equation}
Hence $\varrho(\zeta) = \zeta^2 - 2\zeta + 1 = (\zeta-1)^2$ and
$\sigma(\zeta) = 1$ (the center-point quadrature rule for $f$).

Apply the generating-polynomial test~\eqref{eq:gen-poly-order}.  Expanding
$\varrho(e^h)$ term by term:
\begin{equation}
  \varrho(e^h)
  = e^{2h} - 2e^h + 1
  = h^2 + \frac{h^4}{12} + \frac{h^6}{360} + \cdots
\end{equation}
(all odd powers cancel because $e^{2h} - 2e^h + 1$ is symmetric under $h
\mapsto -h$ after recentering).  With $\sigma(\zeta)=1$:
\begin{equation}
  \label{eq:k2-residual}
  \varrho(e^h) - h^2\sigma(e^h)
  = \underbrace{h^2}_{\text{cancel}} + \frac{h^4}{12} + \frac{h^6}{360} + \cdots
    - h^2
  = \frac{h^4}{12} + \frac{h^6}{360} + \cdots
  = \OO(h^4).
\end{equation}
Comparing with condition~\eqref{eq:order-condition}: $h^{p+2} = h^4 \Rightarrow$
\boxed{p = 2}. The method is \textbf{globally second-order accurate} in $h$.

The leading residual coefficient $\tfrac{1}{12}$ is the \textbf{error constant}
$C = \sigma_2 = \tfrac{1}{12}$ (HNW Exercise~5, §III.10).  It governs the
asymptotic error via (HNW Eq.~10.36--10.37, p.~471):
\begin{equation}
  \label{eq:error-ode}
  \kappa''(x) = \frac{\partial f}{\partial y}\bigl(x,\, y(x)\bigr)\,\kappa(x)
                - \frac{1}{12}\, y^{(4)}(x),
\end{equation}
where $\kappa(x)$ satisfies $y_n - y(x_n) = \kappa(x_n)\,h^2 + \OO(h^3)$.
For orbital mechanics, $y^{(4)}$ is the fourth derivative of the position
trajectory; a smooth orbit (low eccentricity) keeps $y^{(4)}$ small, making
the error constant favorable.

\subsection{The Rounding-Error Trap: Why the $\delta$-Form Is Mandatory}
\label{ssec:double-root}

The polynomial $\varrho(\zeta) = (\zeta-1)^2$ has a \textbf{double root at
$\zeta = 1$}.  HNW Definition~10.1 (p.~467) allows roots on the unit circle
with multiplicity at most 2 --- our method sits exactly on this boundary.

What this means in practice (HNW §III.10, pp.~472--473, Fig.~10.3): a
\emph{direct} implementation that stores $y_0, y_1, \ldots$ and applies the
recurrence $y_{n+1} = 2y_n - y_{n-1} + h^2 f_n$ accumulates rounding errors
that grow without bound as $h \to 0$.  The double root causes the companion
matrix of the difference equation to be no longer power-bounded.

The \textbf{fix} is the $\delta$-increment form (HNW Eq.~10.38$'$, p.~472):
\begin{equation}
  \label{eq:delta-stable}
  \delta_n - \delta_{n-1} = h^2 f_n,
  \qquad
  y_{n+1} - y_n = \delta_n.
\end{equation}
This splits $\varrho(\zeta)$ as $(\zeta-1)(\zeta-1)$ and introduces $\delta_n =
y_{n+1} - y_n$ as a new variable.  The resulting companion matrix is
power-bounded (HNW Eq.~10.30), and rounding errors remain bounded.  This is
exactly the formulation used in \texttt{StormerStep::step()}.

A second pitfall (HNW p.~473, note~b): when computing \emph{starting values}
$\delta_0, \delta_1, \ldots$, do \textbf{not} form the differences
$(y_1 - y_0)/h$, $(y_2 - y_1)/h$ from separately computed RK values --- that
difference is numerically unstable.  Instead, accumulate the increments
directly from the sub-steps of the RK starting procedure, just as
\texttt{StormerStep} does internally.

% ---------------------------------------------------------------------------
\section{Störmer's Modified Midpoint Algorithm}
\label{sec:stormer-alg}
% ---------------------------------------------------------------------------

\paragraph{Notation.} We introduce $\theta_m$ to denote the \textbf{numerical approximation} to the true solution $y(x_m)$ at sub-step $m$, where $x_m = x_0 + mh$. The exact solution is $y(x_m)$; its approximation is $\theta_m$. The algorithm computes a sequence of approximations:
\[
\theta_0 \approx y(x_0),\quad \theta_1 \approx y(x_1),\quad \ldots,\quad \theta_M \approx y(x_0 + H).
\]
The velocity approximation at the final point is denoted $z_M \approx y'(x_0 + H)$.

\paragraph{Why $\theta$ instead of $y$?} In numerical analysis, it is standard to distinguish the \emph{exact} solution $y(x)$ (unknown, continuous) from the \emph{computed} approximation $\theta_m$ (known, discrete). The error at step $m$ is $y(x_m) - \theta_m$. This distinction becomes crucial when analyzing convergence and error propagation.

\bigskip\noindent
The complete algorithm for $M$ sub-steps of size $h = H/M$ is:

\begin{enumerate}
  \item \textbf{Initialize.}
    \begin{align}
      \theta_0 &= y_0, \label{eq:alg-init-pos} \\
      \delta_0 &= h\, v_0 + \tfrac{h^2}{2}\, f(x_0,\, \theta_0). \label{eq:alg-init-delta}
    \end{align}
    Set $\theta_1 = \theta_0 + \delta_0$.

  \item \textbf{Sub-step loop} for $m = 1, \ldots, M-1$:
    \begin{align}
      f_m        &= f(x_0 + m h,\; \theta_m), \label{eq:alg-force} \\
      \delta_m   &= \delta_{m-1} + h^2\, f_m, \label{eq:alg-delta} \\
      \theta_{m+1} &= \theta_m + \delta_m. \label{eq:alg-pos}
    \end{align}

  \item \textbf{Final derivative} (velocity estimate at $x_0 + H$):
    \begin{equation}
      \label{eq:alg-vel}
      z_M = \frac{\delta_{M-1}}{h} + \frac{h}{2}\, f(x_0 + H,\; \theta_M).
    \end{equation}
\end{enumerate}

The output is $(\theta_M,\; z_M) \approx (y(x_0+H),\; y'(x_0+H))$,
the numerical approximations to the true position and velocity at $x_0 + H$,
with local truncation error $\OO(h^2)$.

Total function evaluations of $f$: one per sub-step $m = 0, \ldots, M$, i.e.\
$M+1$ total, but $f_0$ is supplied from the previous step, so \emph{$M$ new
evaluations}.

% ---------------------------------------------------------------------------
\section{Equivalence to Störmer-Verlet / Leapfrog}
% ---------------------------------------------------------------------------

Define the half-integer velocity $v_{m+1/2} \coloneqq \delta_m / h$.  Rewriting the
algorithm in terms of $v$:

\begin{align}
  v_{1/2}       &= \frac{\delta_0}{h} = v_0 + \frac{h}{2}\, f_0
                   && \text{(half-kick)}, \label{eq:lf-halfkick} \\[4pt]
  \theta_{m+1}  &= \theta_m + h\, v_{m+1/2}
                   && \text{(drift)}, \label{eq:lf-drift} \\[4pt]
  v_{m+3/2}     &= v_{m+1/2} + h\, f_{m+1}
                   && \text{(full kick, } m = 0, \ldots, M-2\text{)}, \label{eq:lf-kick}
\end{align}
and the final velocity:
\begin{equation}
  \label{eq:lf-finalvel}
  z_M = v_{M-1/2} + \frac{h}{2}\, f_M.
\end{equation}
Equation~\eqref{eq:lf-finalvel} can be read as: take the half-integer velocity
$v_{M-1/2}$, and kick it forward by a half step to return to integer time $M$.
Equivalently:
\begin{equation}
  z_M = \frac{v_{M-1/2} + v_{M+1/2}}{2},
  \qquad \text{where } v_{M+1/2} = v_{M-1/2} + h\, f_M.
\end{equation}
So $z_M$ is the arithmetic mean of the bracketing half-integer velocities, a
second-order approximation to $y'(x_0 + H)$.

\begin{remark}[Symplecticity]
  The drift–kick–drift structure \eqref{eq:lf-halfkick}--\eqref{eq:lf-kick} is
  the \emph{Störmer-Verlet} integrator, which is \emph{symplectic} for Hamiltonian
  systems $H = \tfrac12 |v|^2 + V(y)$ with $f = -\nabla V$.  It preserves a
  modified Hamiltonian exactly to all orders.  Importantly, it does NOT suffer from
  secular energy drift.
\end{remark}

% ---------------------------------------------------------------------------
\section{Even Error Series: Gragg's Theorem}
% ---------------------------------------------------------------------------

\begin{theorem}[Gragg 1965]
  \label{thm:gragg}
  Let $y(x)$ be sufficiently smooth.  For Störmer's modified midpoint rule applied
  to $y'' = f(x, y)$ with $M$ sub-steps of $h = H/M$ and the starting/ending
  conventions \eqref{eq:delta0}--\eqref{eq:alg-vel}, the global errors admit
  expansions in \emph{even powers} of $h$ only:
  \begin{align}
    y(x_0 + H) - \theta_M &= a_2(H)\, h^2 + a_4(H)\, h^4 + a_6(H)\, h^6 + \cdots,
    \label{eq:even-pos} \\
    y'(x_0 + H) - z_M    &= b_2(H)\, h^2 + b_4(H)\, h^4 + b_6(H)\, h^6 + \cdots,
    \label{eq:even-vel}
  \end{align}
  where the coefficients $a_{2k}(H)$, $b_{2k}(H)$ are independent of $M$.
\end{theorem}

\paragraph{Sketch of proof (time-reversibility argument).}
The Störmer recurrence \eqref{eq:alg-init-pos}--\eqref{eq:alg-vel} is
\emph{time-reversible}: reversing $h \to -h$ and the sequence of $\theta_m$
recovers the original sequence.  A time-reversible method cannot accumulate odd
powers of $h$ in its global error, because odd-power terms change sign under
$h \to -h$ while the exact solution does not.  More formally, one expands
$\theta_m(h)$ as a power series in $h$ and shows that all odd-power terms cancel
identically (see Hairer, Nørsett \& Wanner, Vol.~I, §II.8 for the complete
argument).

\begin{corollary}[Richardson extrapolation efficiency]
  Because the error series \eqref{eq:even-pos} contains only $h^2$, $h^4$, \ldots,
  a single Richardson extrapolation step using results at $h$ and $h/2$ eliminates
  the $h^2$ term and yields $\OO(h^4)$ accuracy.  In general, $k$ extrapolation
  steps yield $\OO(h^{2k+2})$ accuracy from $\OO(h^2)$ base integrations.  This
  is the foundation of the Bulirsch-Stoer method.
\end{corollary}

% ---------------------------------------------------------------------------
\section{Richardson Extrapolation Setup (Bulirsch-Stoer)}
% ---------------------------------------------------------------------------

Evaluate the modified midpoint rule with sub-step counts $M_1 < M_2 < M_3 < \cdots$
(the Bulirsch-Stoer sequence: $\{2, 4, 6, 8, 12, 16, \ldots\}$), giving
approximations $T_{j,0} \coloneqq \theta_{M_j}(H)$.

Define the extrapolation table recursively:
\begin{equation}
  T_{j,k} = T_{j,k-1} + \frac{T_{j,k-1} - T_{j-1,k-1}}
                              {\bigl(M_j / M_{j-k}\bigr)^2 - 1},
  \qquad k = 1, 2, \ldots, j.
\end{equation}
Each extrapolation step eliminates the next order in $h^2$.  The diagonal $T_{k,k}$
converges to $y(x_0 + H)$ as $k$ increases.

\begin{remark}[Computational cost vs.\ accuracy]
  For the first-order system $[y', f(x,y)]^T$ of dimension $2n$, the modified
  midpoint rule needs $M$ evaluations of $f$ per column.  Störmer's rule needs only
  $M/2$ evaluations per column (the NR3 code uses $n_\text{eff} = M/2$ leapfrog
  steps of $h_\text{eff} = 2h = 2H/M$, as shown below) for the same total interval
  $H$.  Since $f$ can be expensive (e.g., gravitational force with many bodies),
  this factor of two is significant.
\end{remark}

% ---------------------------------------------------------------------------
\section{Analysis of the NR3 Implementation}
% ---------------------------------------------------------------------------

\subsection{Parameterization}

NR3's \texttt{StepperStoerm::dy()} uses:
\begin{itemize}
  \item \texttt{nstep} $= M$ (from Bulirsch-Stoer sequence $\{2,4,6,8,\ldots\}$).
  \item \texttt{h} $= H/M$ (unit sub-step).
  \item \texttt{h2} $= 2h = 2H/M$ (actual loop step, double the unit sub-step).
  \item \texttt{nstp2} $= M/2$ (number of loop iterations).
\end{itemize}

\subsection{Algorithm in NR3 variables}

\textbf{Initialization:}
\begin{align*}
  \text{ytemp}[i]   &\leftarrow y[i]                && \theta_0 = y_0, \\
  \text{ytemp}[n+i] &\leftarrow y'[i] + h\cdot f_0  && v^* = v_0 + h\,f_0
                                                       = v_{1/2} \cdot
                                                       \underbrace{2}_{h_\text{eff}/h}.
\end{align*}
\noindent In other words, the stored ``velocity'' \texttt{ytemp}[$n+i$] is the
half-integer velocity $v_{1/2} = v_0 + (h/2)\,f_0$ \emph{scaled by 2}:
\texttt{ytemp}[$n+i$] $= 2\,v_{1/2}$.  Then the drift step below with factor
$h_\text{eff} = 2h$ is:
\[
  \theta_1 = \theta_0 + h_\text{eff} \cdot \text{ytemp}[n+i]/2
           = \theta_0 + h_\text{eff}\,v_{1/2}
           = \theta_0 + 2h\left(v_0 + \tfrac{h}{2}\,f_0\right)
           = \theta_0 + 2h\,v_0 + h^2\,f_0.
\]
Alternatively, viewing the NR3 scheme directly as Störmer-Verlet with step
$h_\text{eff} = 2h$:

\begin{equation}
  \label{eq:nr3-leapfrog}
  \text{Loop } m = 1, \ldots, n_\text{eff}:\quad
  \begin{cases}
    \theta_m = \theta_{m-1} + h_\text{eff}\,v_{m-1/2}, \\
    v_{m+1/2} = v_{m-1/2} + h_\text{eff}\,f(x_0 + m\,h_\text{eff},\;\theta_m),
      & m < n_\text{eff}.
  \end{cases}
\end{equation}

\noindent Final velocity:
$z = v_{n_\text{eff}-1/2} + (h_\text{eff}/2)\,f(x_0 + H, \theta_{n_\text{eff}})$.

\subsection{Cost comparison}

\begin{center}
\begin{tabular}{@{}lccc@{}}
  \toprule
  Method & State dimension & Evals of $f$ per column & Error order \\
  \midrule
  Midpoint (first-order 2n-system) & $2n$ & $M$   & $\OO(h^2)$ \\
  Störmer (NR3)                    & $n$  & $M/2$ & $\OO(h_\text{eff}^2)$ \\
  \bottomrule
\end{tabular}
\end{center}

Both methods embed in a Bulirsch-Stoer extrapolation with an even-power error
series and converge to the same accuracy.  Störmer uses half the evaluations of $f$
at the cost of doubling the effective sub-step size, which the extrapolation corrects.

% ---------------------------------------------------------------------------
\section{Is Inheritance Necessary?}
% ---------------------------------------------------------------------------

NR3's design uses \texttt{struct StepperStoerm : StepperBS<D>} \emph{solely} for
code reuse of the Bulirsch-Stoer extrapolation machinery (\texttt{odeint} driver,
column sequence, Richardson table, stepsize control).  The method \texttt{dy()} is
entirely independent of the base class algorithms; it overrides only the
\emph{sub-step integration} while reusing everything else.

\paragraph{From first principles, the minimal reusable components are:}
\begin{enumerate}
  \item \textbf{\texttt{StormerStep}}: executes $n_\text{eff}$ leapfrog sub-steps
        over a given interval $H$.  No base class required.  Template on
        \texttt{DerivativeType} and \texttt{ContainerT}.
  \item \textbf{Error estimation}: step-doubling (compare $n_\text{eff}$ steps vs.\
        $2\,n_\text{eff}$ steps) or Richardson extrapolation as a separate,
        composable component.
  \item \textbf{Step-size control}: compose with the existing
        \texttt{ComputePIStepSize} / \texttt{PIStepSizeControl} classes already
        present in the library.
\end{enumerate}

Composition over inheritance, exactly consistent with how
\texttt{CalculateNewYAndError} and \texttt{IntegrateWithPIControl} were designed.

% ---------------------------------------------------------------------------
\section{Harmonic Oscillator Test Case}
% ---------------------------------------------------------------------------

The standard test for Störmer's rule is the harmonic oscillator:
\begin{equation}
  y'' = -\omega^2 y, \qquad y(0) = 1,\quad y'(0) = 0,\quad \omega = 1.
\end{equation}
Exact solution: $y(t) = \cos t$, $y'(t) = -\sin t$, with conserved energy:
\begin{equation}
  E(t) = \tfrac{1}{2}\bigl(y'(t)\bigr)^2 + \tfrac{1}{2}\bigl(y(t)\bigr)^2 = \tfrac12.
\end{equation}

Störmer-Verlet (symplectic) conserves a modified energy exactly, so the
energy error remains bounded for all time (no secular drift), unlike non-symplectic
methods.

\paragraph{Expected convergence.}
For a single step of size $H = 0.5$ with $n_\text{eff}$ leapfrog sub-steps of
$h_\text{eff} = H/n_\text{eff}$:
\begin{align*}
  |y(H) - \theta_M| &\approx \tfrac{H^3}{24}\,|y'''(0)|\cdot h_\text{eff}^2
                             = \tfrac{H^3 \omega^3}{24}\cdot h_\text{eff}^2
                             = \OO(h_\text{eff}^2).
\end{align*}
Doubling $n_\text{eff}$ (halving $h_\text{eff}$) should reduce the error by a
factor of 4.  This is the convergence test in the unit tests.

% ---------------------------------------------------------------------------
\section{Numerov's Method: The Implicit Upgrade}
\label{sec:numerov}
% ---------------------------------------------------------------------------

Section~\ref{sec:order} showed that the three-point stencil~\eqref{eq:stormer-3pt}
achieves order 2 because $\varrho(e^h) - h^2\sigma(e^h) = \OO(h^4)$ with
$\sigma(\zeta) = 1$.  The same stencil on the left-hand side, but a richer
choice of $\sigma$, can reach $\OO(h^6)$ --- i.e.\ order 4 --- with \emph{no
extra history}.

\subsection{Notation Bridge: HNW Polynomials vs NR3/Our Variables}
\label{ssec:notation-bridge}

The analysis in \S\ref{sec:order} introduced Hairer--N{\o}rsett--Wanner
(HNW) notation.  Here is a translation table bridging their polynomial
formalism with the discrete variables used earlier in this document and in
NR3:

\begin{center}
\begin{tabular}{@{}lll@{}}
\toprule
\textbf{Symbol} & \textbf{Meaning} & \textbf{Where it appears} \\
\midrule
$\varrho(\zeta)$ & First characteristic polynomial & HNW Eq.~10.20, p.~467 \\
                & $\sum_{i=0}^k \alpha_i \zeta^i$ & (stability) \\
$\sigma(\zeta)$ & Second characteristic polynomial & HNW Eq.~10.20, p.~467 \\
                & $\sum_{i=0}^k \beta_i \zeta^i$ & (accuracy) \\
$\zeta$         & Formal variable; $E = \zeta$ is the & HNW \S III.1 \\
                & shift operator: $E y_n = y_{n+1}$ & \\
$\alpha_i$      & Coefficients of LHS stencil & Eq.~\eqref{eq:general-multistep} \\
                & (position weights) & \\
$\beta_i$       & Coefficients of RHS quadrature & Eq.~\eqref{eq:general-multistep} \\
                & (force weights) & \\
$k$             & Step number (history length $k$) & HNW \S III.10 \\
                & Method uses $y_{n},\dots,y_{n+k}$ & \\
$h$             & Step size & Same throughout \\
$H$             & Total interval (macro-step) & \S\ref{sec:stormer-alg} \\
$y_n$           & Exact solution at $x_n = x_0 + nh$ & Continuous $y(x_n)$ \\
$\theta_n$      & \textbf{Computed} approximation to $y_n$ & \S\ref{sec:stormer-alg} \\
$f_n$           & Force $f(x_n, y_n)$ or $f(x_n, \theta_n)$ & \S\ref{sec:stormer-alg} \\
$\delta_n$      & Position increment $\theta_{n+1}-\theta_n$ & Eq.~\eqref{eq:delta-recurrence} \\
$v_{1/2}$       & Half-integer velocity (leapfrog) & Eq.~\eqref{eq:lf-halfkick} \\
\bottomrule
\end{tabular}
\end{center}

\paragraph{Connection to our earlier notation.}
In \S\ref{sec:stormer-3pt}--\ref{sec:stormer-alg} we wrote the method as
\[
  \theta_{m+1} - 2\theta_m + \theta_{m-1} = h^2 f_m.
\]
In HNW polynomial form this is
\[
  \varrho(\zeta) = \zeta^2 - 2\zeta + 1 = (\zeta-1)^2,
  \qquad
  \sigma(\zeta) = 1.
\]
The coefficients are $(\alpha_0,\alpha_1,\alpha_2) = (1,-2,1)$ and
$(\beta_0,\beta_1,\beta_2) = (0,1,0)$.  The shift operator $E$ acts as
$E\theta_m = \theta_{m+1}$, so the compact operator form is
$\varrho(E)\theta_m = h^2 \sigma(E) f_m$.

\paragraph{NR3 notation.}
NR3 uses `nstep` for the number of sub-steps $M$, `h` for $H/M$, `h2`
for $2h$, and `nstp2` for $M/2$.  Their `ytemp[n+i]` stores the
half-integer velocity $v_{m+1/2}$ scaled by 2 (see
\S\ref{sec:nr3-analysis}).

\subsection{Derivation via the Order Condition}

Keep $\varrho(\zeta) = \zeta^2 - 2\zeta + 1$ and seek $\sigma(\zeta) =
\beta_0 + \beta_1\zeta + \beta_2\zeta^2$ (symmetric: $\beta_0 = \beta_2$) such
that $\varrho(e^h) - h^2\sigma(e^h) = \OO(h^6)$.

From~\eqref{eq:k2-residual}, $\varrho(e^h) = h^2 + \tfrac{h^4}{12} +
\tfrac{h^6}{360} + \cdots$.  We need $h^2\sigma(e^h)$ to match through the
$h^5$ term.  Expanding with $\beta_0 = \beta_2 \equiv \beta$,
$\beta_1 \equiv \gamma$:
\[
  h^2\sigma(e^h)
  = h^2\bigl(\beta e^{2h} + \gamma e^h + \beta\bigr)
  = h^2\bigl[(2\beta + \gamma) + 2h\beta + h^2(2\beta + \tfrac\gamma2) + \cdots\bigr].
\]
Matching through $h^5$ against $\varrho(e^h)$ forces $2\beta + \gamma = 1$
(the $h^2$ term), and solving the system for the $h^4$ match gives
$\beta = \tfrac{1}{12}$, $\gamma = \tfrac{10}{12}$.  Hence:
\begin{equation}
  \label{eq:numerov-sigma}
  \sigma(\zeta) = \frac{\zeta^2 + 10\zeta + 1}{12}.
\end{equation}
The resulting method is (HNW Eq.~10.12, p.~464; B.~Numerov 1924):
\begin{equation}
  \label{eq:numerov}
  \boxed{y_{n+1} - 2y_n + y_{n-1}
  = \frac{h^2}{12}\Bigl(f_{n+1} + 10\,f_n + f_{n-1}\Bigr),}
\end{equation}
called \textbf{Numerov's method}.  The only difference from~\eqref{eq:stormer-3pt}
is the $f_{n+1}$ term on the right: one additional force evaluation at the new
position.

\subsection{Order Verification}

Apply the generating-polynomial test~\eqref{eq:gen-poly-order}:
\begin{align}
  h^2\sigma(e^h)
  &= \frac{h^2}{12}\bigl(e^{2h} + 10e^h + 1\bigr) \notag \\
  &= h^2 + \frac{h^4}{12} + \frac{h^6}{144} + \cdots
     \label{eq:numerov-sigma-expand}
\end{align}
and recall $\varrho(e^h) = h^2 + \tfrac{h^4}{12} + \tfrac{h^6}{360} + \cdots$.
Subtracting:
\begin{equation}
  \label{eq:numerov-residual}
  \varrho(e^h) - h^2\sigma(e^h)
  = \Bigl(\frac{1}{360} - \frac{1}{144}\Bigr)h^6 + \OO(h^8)
  = -\frac{1}{240}\,h^6 + \OO(h^8)
  = \OO(h^6).
\end{equation}
Comparing with~\eqref{eq:order-condition}: $h^{p+2} = h^6 \Rightarrow$
\boxed{p = 4}.  \textbf{Numerov's method is globally fourth-order accurate.}

The error constant is $C = \sigma^*_3 = 0$ (Table~10.2, HNW p.~464) --- the
leading error is $-h^6/240 \cdot y^{(6)}$ per step, contributing:
\begin{equation}
  \label{eq:numerov-error}
  \kappa''(x) = \frac{\partial f}{\partial y}\,\kappa - \frac{1}{240}\,y^{(6)}(x)
  \qquad\bigl(\text{global error: } y_n - y(x_n) = \kappa(x_n)\,h^4 + \OO(h^5)\bigr).
\end{equation}

\subsection{The Order-Barrier Context}

HNW Theorem~10.4 (p.~468) states that for a stable $k$-step Störmer method:
\begin{equation}
  p \leq k + 2 \;\text{ if } k \text{ is even},
  \qquad
  p \leq k + 1 \;\text{ if } k \text{ is odd}.
\end{equation}
For $k = 2$ (our stencil): $p_\text{max} = 4$.  Numerov saturates this bound ---
it is the \emph{highest-order stable method} on the three-point stencil.  The
explicit method (order 2) leaves two orders on the table.

\subsection{PECECE Implementation}
\label{ssec:pecece}

Numerov's method is implicit: $y_{n+1}$ appears on both sides via $f_{n+1} =
f(x_{n+1}, y_{n+1})$.  For nonlinear $f$ (e.g.\ gravitational force), we solve
using \textbf{PECECE} (Predict--Evaluate--Correct--Evaluate--Correct--Evaluate):

\begin{enumerate}
  \item \textbf{Predict.}  Use the explicit Störmer formula as first guess:
    \begin{equation}
      \label{eq:pece-predict}
      y_{n+1}^{(P)} = 2y_n - y_{n-1} + h^2 f_n.
    \end{equation}
  \item \textbf{Evaluate.}  $f_{n+1}^{(P)} = f\!\left(x_{n+1},\, y_{n+1}^{(P)}\right)$.
  \item \textbf{Correct.}
    \begin{equation}
      \label{eq:pece-correct1}
      y_{n+1}^{(1)} = 2y_n - y_{n-1}
                    + \frac{h^2}{12}\Bigl(f_{n+1}^{(P)} + 10 f_n + f_{n-1}\Bigr).
    \end{equation}
  \item \textbf{Evaluate.}  $f_{n+1}^{(1)} = f\!\left(x_{n+1},\, y_{n+1}^{(1)}\right)$.
  \item \textbf{Correct again.}
    \begin{equation}
      \label{eq:pece-correct2}
      y_{n+1} = 2y_n - y_{n-1}
              + \frac{h^2}{12}\Bigl(f_{n+1}^{(1)} + 10 f_n + f_{n-1}\Bigr).
    \end{equation}
  \item \textbf{Evaluate.}  $f_{n+1} = f\!\left(x_{n+1},\, y_{n+1}\right)$
        \quad (stored for next step).
\end{enumerate}

\paragraph{Velocity.}
Numerov delivers positions.  Velocity at integer time is recovered via the
leapfrog formula (order 2 in $h$):
\begin{equation}
  \label{eq:numerov-velocity}
  v_{n+1} = \frac{y_{n+1} - y_n}{h} + \frac{h}{2}\,f_{n+1}.
\end{equation}
\emph{Position accuracy is $\OO(h^4)$; velocity accuracy is $\OO(h^2)$.}  For
applications where velocity precision matters as much as position (e.g.\
orbit determination), the velocity can be improved via central differences:
$v_n \approx (y_{n+1} - y_{n-1})/(2h) + \OO(h^2)$ --- same order.

\paragraph{Order of PECECE.}
With an order-2 predictor and order-4 corrector, PECE alone is order~3: the
predictor error $\OO(h^2)$ contaminates $f_{n+1}^{(P)}$ by $\OO(h^2)$, adding
$\tfrac{h^2}{12}\cdot\OO(h^2) = \OO(h^4)$ per step to the corrector residual,
giving $\OO(h^3)$ global error.  The second correction (PECECE) reduces the
contamination to $\OO(h^6)$ per step --- below the corrector's own leading
error of $\OO(h^6)$ --- so the full order-4 accuracy is recovered.

\subsection{Cost Comparison}

\begin{center}
\begin{tabular}{@{}lcccc@{}}
  \toprule
  Method & Order $p$ & $f$ evals/step & Error constant $C$ \\
  \midrule
  Explicit Störmer ($k=2$) & 2 & 1 & $\sigma_2 = \tfrac{1}{12}$ \\
  Numerov PECECE ($k=2$)   & 4 & 3 & $-\tfrac{1}{240}$\\
  \bottomrule
\end{tabular}
\end{center}

Halving $h$ reduces error by $2^2 = 4$ for explicit Störmer vs.\ $2^4 = 16$
for Numerov --- a 4$\times$ gain per step halving.  The cost of 3 $f$-evals vs.\
1 is recovered after the first tolerance tightening: to reach the same accuracy,
Numerov needs $\approx 2.5\times$ fewer steps than explicit Störmer (since
$4^{1/4} / 2^{1/2} \approx 1.41$ ratio in step size for equal error).

\subsection{Starting Procedure}

Numerov is a two-step method: it needs $y_0$ and $y_1$ to start.  Given initial
conditions $(y_0, v_0)$, one accurate starting step is required.
Theorem~\ref{thm:convergence} demands $y(x_1) - y_1 = \OO(h^5)$ for the
method to retain its full order-4 convergence.  Options:

\begin{itemize}
  \item \textbf{Taylor expansion:}
    $y_1 = y_0 + hv_0 + \tfrac{h^2}{2}f_0 + \tfrac{h^3}{6}f_0' + \tfrac{h^4}{24}f_0''$
    (requires derivative information, gives $\OO(h^5)$, satisfies requirement).
  \item \textbf{Runge-Kutta-Nyström order~4:} one step of a 4th-order Nyström
    method (HNW Table~14.2, p.~285) gives $\OO(h^5)$ accuracy, at the cost of 3--4
    $f$-evaluations for the single startup step.
  \item \textbf{Two steps of explicit Störmer} (order~2): gives only $\OO(h^3)$,
    limiting global convergence to $\OO(h^2)$.  Sufficient only for testing.
\end{itemize}

\texttt{NumerovStep::startup()} uses one step of the leapfrog (Störmer-Verlet)
with $n_\text{eff}$ sub-steps as a high-accuracy starter: for large
$n_\text{eff}$, this approximates the exact solution to whatever precision is
needed.  In production, supply a 4th-order starter.

% ---------------------------------------------------------------------------
\section{Astrodynamics Context: Cowell's Method, Encke's Method, and Perturbations}
\label{sec:astrodynamics}
% ---------------------------------------------------------------------------

In orbital mechanics, the equations of motion for a satellite under perturbations are:
\begin{equation}
\ddot{\mathbf{r}} = -\frac{\mu}{r^3}\mathbf{r} + \mathbf{a}_p
\end{equation}
where $\mathbf{a}_p$ is the sum of all perturbing accelerations. Two main approaches exist for numerical integration:

\subsection{Cowell's Method (Direct Integration)}
\label{ssec:cowell}

Cowell's method \cite{Bate1971,Curtis2005,Hintz2015} directly integrates the total acceleration:
\begin{align}
\dot{\mathbf{r}} &= \mathbf{v} \\
\dot{\mathbf{v}} &= -\frac{\mu}{r^3}\mathbf{r} + \mathbf{a}_p
\end{align}

\textbf{Advantages:}
\begin{itemize}
\item Simplicity of formulation and implementation
\item Any number of perturbations can be handled simultaneously
\item Obtains $\mathbf{r}$ and $\mathbf{v}$ directly
\end{itemize}

\textbf{Disadvantages:}
\begin{itemize}
\item Near large attracting bodies, smaller steps are required (affects time and roundoff error)
\item Approximately 10$\times$ slower than Encke's method \cite{Bate1971}
\item Roundoff error accumulates rapidly in double-precision for interplanetary flights with small step sizes
\end{itemize}

\subsection{Encke's Method (Osculating Orbit)}
\label{ssec:encke}

Encke's method \cite{Bate1971,Hintz2015} integrates the \textit{deviation} from a reference (osculating) two-body orbit rather than the total state. Let $\mathbf{\rho}(t)$ be the osculating orbit position (the path the satellite would follow if perturbations vanished at epoch $t_0$). The deviation $\delta\mathbf{r} = \mathbf{r} - \mathbf{\rho}$ satisfies:
\begin{equation}
\delta\ddot{\mathbf{r}} = \mathbf{a}_p + \frac{\mu}{\rho^3}\left[(1 - 2q)^{-3/2}\mathbf{r} - \delta\mathbf{r}\right]
\end{equation}
where $q = \delta\mathbf{r} \cdot (2\mathbf{r} - \delta\mathbf{r}) / r^2$.

\textbf{Advantages:}
\begin{itemize}
\item Faster than Cowell's method (deviations change slowly)
\item Better accuracy for small perturbations
\item Less accumulation of roundoff error
\end{itemize}

\textbf{Disadvantages:}
\begin{itemize}
\item More complex formulation
\item Requires \textit{rectification} (resetting the osculating orbit) when $\delta r/r$ exceeds a tolerance
\item The term $(1 - 2q)^{-3/2}$ requires careful numerical handling (difference of nearly equal quantities)
\end{itemize}

\subsection{Störmer/Numerov vs. Cowell for Orbital Mechanics}
\label{ssec:stormer-vs-cowell}

Our \texttt{StörmerStep} and \texttt{NumerovStep} implementations are \textbf{integration methods}, while Cowell's and Encke's are \textbf{formulation frameworks}. The relationship:

\begin{center}
\begin{tabular}{lll}
\toprule
\textbf{Formulation} & \textbf{Integration Method} & \textbf{Use Case} \\
\midrule
Cowell's & Störmer (order 2) or Numerov (order 4) & Direct integration of $y'' = f(x,y)$ \\
Cowell's & RK4, Adams-Moulton & General first-order reduction \\
Encke's & Störmer/Numerov & Deviation from osculating orbit \\
\bottomrule
\end{tabular}
\end{center}

\textbf{Why Numerov outperforms standard Cowell integration:}
\begin{enumerate}
\item \textbf{Order 4 accuracy} vs. order 2 for explicit Störmer or standard RK4 on first-order form
\item \textbf{Same stencil} (3 points) as explicit Störmer, but implicit correction achieves higher order
\item \textbf{Symplectic-like} energy behavior for near-Keplerian orbits (nearly circular)
\item \textbf{Fewer function evaluations} than equivalent RK4: 3 $f$-evals/step vs. 4 for RK4
\end{enumerate}

Bate et al.\ \cite{Bate1971} note that for orbital problems, the \textbf{Gauss-Jackson} (sum-squared, $\Sigma^2$) method is "one of the best, and most used, numerical integration methods for trajectory problems of the Cowell and Encke type." The Gauss-Jackson method is also a multistep predictor-corrector for second-order equations, similar in spirit to Numerov but with higher order and specialized for orbital mechanics.

\subsection{Common Perturbations in Astrodynamics}
\label{ssec:perturbations}

Based on Curtis \cite{Curtis2005}, Hintz \cite{Hintz2015}, and Bate \cite{Bate1971}, the primary perturbations affecting Earth satellites are:

\begin{enumerate}
\item \textbf{Earth Oblateness ($J_2$)}: The Earth's equatorial bulge causes secular variations in $\Omega$ (RAAN) and $\omega$ (argument of perigee). The nodal regression rate is:
\begin{equation}
\dot{\Omega} = -\frac{3J_2}{2}\sqrt{\frac{\mu}{a^3}}\left(\frac{R_\oplus}{a(1-e^2)}\right)^2\cos i
\end{equation}

\item \textbf{Atmospheric Drag}: Significant for LEO satellites. The acceleration is:
\begin{equation}
\mathbf{a}_d = -\frac{1}{2}\rho v_{rel}^2 \frac{C_D A}{m}\frac{\mathbf{v}_{rel}}{|\mathbf{v}_{rel}|}
\end{equation}
where $\rho$ is atmospheric density, $C_D \approx 2.0$--$2.2$ for typical satellites.

\item \textbf{Third-Body Gravity} (Moon, Sun): The lunar perturbation acceleration on an Earth satellite is:
\begin{equation}
\mathbf{a}_p = \mu_M\left(\frac{\mathbf{r}_{ms}}{r_{ms}^3} - \frac{\mathbf{r}_{Me}}{r_{Me}^3}\right)
\end{equation}
where $\mathbf{r}_{ms}$ is the vector from Moon to satellite, $\mathbf{r}_{Me}$ from Moon to Earth.

\item \textbf{Solar Radiation Pressure}: For spacecraft with large area-to-mass ratio:
\begin{equation}
\mathbf{a}_{SRP} = -p_{SR}\frac{A}{m}\cos\theta\,\mathbf{\hat{s}}
\end{equation}
where $p_{SR} \approx 4.5 \times 10^{-6}$ N/m$^2$ at 1 AU.
\end{enumerate}

\textbf{Special vs. General Perturbations:}
\begin{itemize}
\item \textbf{Special perturbations}: Direct numerical integration (Cowell, Encke) --- computes specific trajectories for given initial conditions.
\item \textbf{General perturbations}: Analytic integration of series expansions (variation of parameters, Lagrange planetary equations) --- provides analytical insight but is more complex.
\end{itemize}

% ---------------------------------------------------------------------------
\subsection{Application: Numerov Orbit Propagation Demo}
\label{ssec:numerov-demo}
% ---------------------------------------------------------------------------

The \texttt{NumerovOrbitDemo.cpp} implementation demonstrates two scenarios:

\paragraph{Toy Example: Pure Two-Body LEO (400~km, 28.5° inclination)}
\begin{itemize}
\item \textbf{Purpose}: Verify Numerov implementation correctness
\item \textbf{Physics}: Central gravity only ($\mathbf{a} = -\mu\mathbf{r}/r^3$)
\item \textbf{Validation}: Energy conservation to $\sim$0.04\% over 10 orbits
\item \textbf{Expected}: Perfectly periodic orbit with constant energy
\end{itemize}

\paragraph{Real-World Example: LEO with $J_2$ Perturbation}
\begin{itemize}
\item \textbf{Purpose}: Demonstrate dominant LEO perturbation effect
\item \textbf{Physics}: Two-body + Earth oblateness ($J_2 = 1.08263 \times 10^{-3}$)
\item \textbf{Perturbation formula} (Curtis \S10.5):
\begin{equation}
\mathbf{a}_{J_2} = -\frac{3\mu J_2 R_\oplus^2}{2r^5}
\begin{bmatrix}
x(1 - 5z^2/r^2) \\
y(1 - 5z^2/r^2) \\
z(3 - 5z^2/r^2)
\end{bmatrix}
\end{equation}
\item \textbf{Observable effects}:
  \begin{itemize}
  \item Nodal regression: $-7.08°$/day (Curtis Eq.~4.52)
  \item Perigee advance: $+11.52°$/day (Curtis Eq.~4.53)
  \end{itemize}
\item \textbf{Why nontrivial}: $J_2$ is the dominant LEO perturbation ($\sim$90\% of all perturbation effects)
\end{itemize}

\noindent
\textbf{Build and run} (from \texttt{Cosmos/Source}):
\begin{verbatim}
mkdir -p build && cd build
cmake .. && make NumerovOrbitDemo
./Examples/NumerovOrbitDemo
\end{verbatim}

\noindent
\textbf{Output}: \texttt{numerov\_twobody.csv}, \texttt{numerov\_j2.csv} with columns \texttt{t,orbit,r\_x,r\_y,r\_z,energy,error}.

% ---------------------------------------------------------------------------
\section{Summary}
% ---------------------------------------------------------------------------

\begin{enumerate}
  \item Störmer's rule exploits $y'' = f(x, y)$ (no $y'$ in $f$) via the
        three-point recurrence $y_{m+1} = 2y_m - y_{m-1} + h^2 f_m$.

  \item In $\delta$-form: $\delta_m = \delta_{m-1} + h^2 f_m$, with starting value
        $\delta_0 = hv_0 + \tfrac{h^2}{2}f_0$.  This form is mandatory to
        avoid numerical instability from the double root $\varrho(\zeta) =
        (\zeta-1)^2$ (Section~\ref{ssec:double-root}).

  \item The algorithm is equivalent to Störmer-Verlet (leapfrog), which is
        symplectic for Hamiltonian systems.

  \item ``Order $p$'' means global error $\OO(h^p)$ in step size $h$ (not
        polynomial degree).  The explicit method is order 2;
        Numerov's implicit method~\eqref{eq:numerov} saturates the order barrier
        for the two-step stencil at order 4.

  \item Gragg's theorem guarantees an even-power error series in $h$, enabling
        efficient Bulirsch-Stoer Richardson extrapolation.

  \item NR3's \texttt{StepperStoerm} implements the $k=2$ explicit Störmer in
        one-step extrapolation mode (HNW §II.14, Eq.~14.32) --- the simplest
        special case of the general multistep family (HNW §III.10).

  \item Numerov PECECE (Section~\ref{sec:numerov}) achieves order~4 on the same
        three-point stencil at 3 $f$-evaluations/step vs.\ 1 for explicit
        Störmer.  Implemented as \texttt{NumerovStep}.
\end{enumerate}

\bigskip\noindent
\textbf{References.}

\begin{thebibliography}{99}

\bibitem{HNW} E. Hairer, S. P. N{\o}rsett, G. Wanner,
        \textit{Solving Ordinary Differential Equations~I: Nonstiff Problems},
        2nd ed.\ (Springer, 1993).
        Key sections: §II.14 (Nyström methods, extrapolation for $y''=f(x,y)$),
        §III.10 (multistep Störmer methods, pp.~461--473).
  \bibitem{Gragg1965} W. B. Gragg, ``On extrapolation algorithms for ordinary initial value
        problems,'' \textit{SIAM J.\ Numer.\ Anal.}\ \textbf{2} (1965) 384--403.
\bibitem{Bate1971} R. R. Bate, D. D. Mueller, J. E. White, \textit{Fundamentals of Astrodynamics} (Dover, 1971). \S9.2 (Cowell's Method), \S9.3 (Encke's Method), \S9.6 (Numerical Integration Methods including Gauss-Jackson).
\bibitem{Curtis2005} H. D. Curtis, \textit{Orbital Mechanics for Engineering Students}, 2nd ed.\ (Elsevier, 2005). Ch.~10 (Orbital Perturbations): \S10.2 (Cowell), \S10.3 (Encke), \S10.4 (Atmospheric Drag), \S10.5 (Gravitational Perturbations).
\bibitem{Hintz2015} G. R. Hintz, \textit{Orbital Mechanics and Astrodynamics} (Springer, 2015). Ch.~5: \S5.2 (Cowell's Method), \S5.3 (Encke's Method), \S5.5 (Integration Schemes), \S5.7 (Analytic Formulation of Perturbations).

\bibitem{Vallado2013} D. A. Vallado, \textit{Fundamentals of Astrodynamics and Applications},
        4th ed.\ (Microcosm Press / Springer, 2013).
        Ch.~8 (Perturbation Methods): \S8.6 (Non-spherical Earth), \S8.7 ($J_2$ Secular Effects).

\bibitem{Kaula1966} W. M. Kaula, \textit{Theory of Satellite Geodesy} (Blaisdell, 1966;
        Dover reprint 2000). Ch.~2: spherical harmonic expansion of the geopotential;
        Ch.~3: orbital perturbations.

\bibitem{Heiskanen1967} W. A. Heiskanen and H. Moritz,
        \textit{Physical Geodesy} (Freeman, 1967). \S1-19: Legendre functions and the
        external gravitational potential.

\end{thebibliography}

% ===========================================================================
\section{Gravitational Potential and Spherical Harmonics Expansion}
\label{sec:gravity-harmonics}
% ===========================================================================

This section derives the full spherical harmonic expansion of Earth's gravitational
potential from first principles, obtains the explicit $J_2$ and $J_3$ perturbation
terms, and provides the Cartesian acceleration expressions used in orbit propagators
such as Cosmos GNC.

% ---------------------------------------------------------------------------
\subsection{Coordinate Conventions}
\label{ssec:coords}
% ---------------------------------------------------------------------------

We adopt the standard right-handed \emph{spherical coordinates}
$(r,\theta,\phi)$, where
\begin{itemize}
  \item $r \geq 0$: radial distance from the Earth's centre of mass,
  \item $\theta \in [0,\pi]$: \emph{colatitude} (polar angle measured from
        the north pole along the rotation axis $+z$),
  \item $\phi \in [0, 2\pi)$: \emph{east longitude}.
\end{itemize}
These relate to Cartesian coordinates $(x,y,z)$ --- with $z$ pointing north
along Earth's spin axis --- by
\begin{align}
  x &= r\sin\theta\cos\phi, \label{eq:sph2cart-x}\\
  y &= r\sin\theta\sin\phi, \label{eq:sph2cart-y}\\
  z &= r\cos\theta.         \label{eq:sph2cart-z}
\end{align}
Inverting:
\begin{equation}
  r = \sqrt{x^2+y^2+z^2},\qquad
  \theta = \arccos\!\frac{z}{r},\qquad
  \phi = \mathop{\mathrm{atan2}}(y,x).
  \label{eq:cart2sph}
\end{equation}

\begin{remark}[Latitude vs.\ colatitude]
  Geodetic \emph{latitude} $\varphi$ satisfies $\varphi = \tfrac{\pi}{2}-\theta$,
  so $\cos\theta = \sin\varphi$ and $\sin\theta = \cos\varphi$.  Some references
  (Vallado~\cite{Vallado2013}, Kaula~\cite{Kaula1966}) use latitude; we use
  colatitude throughout for cleaner Legendre notation.
\end{remark}

% ---------------------------------------------------------------------------
\subsection{Transformation Matrix: Spherical to Cartesian}
\label{ssec:transform}
% ---------------------------------------------------------------------------

The spherical orthonormal basis at $(r,\theta,\phi)$ is
\begin{align}
  \hat{e}_r      &= \sin\theta\cos\phi\,\hat{x}
                  + \sin\theta\sin\phi\,\hat{y}
                  + \cos\theta\,\hat{z},
                  \label{eq:e-r}\\
  \hat{e}_\theta &= \cos\theta\cos\phi\,\hat{x}
                  + \cos\theta\sin\phi\,\hat{y}
                  - \sin\theta\,\hat{z},
                  \label{eq:e-theta}\\
  \hat{e}_\phi   &= -\sin\phi\,\hat{x}
                  + \cos\phi\,\hat{y}.
                  \label{eq:e-phi}
\end{align}
These three directions are mutually orthogonal and right-handed:
$\hat{e}_r \times \hat{e}_\theta = \hat{e}_\phi$.
The matrix whose \emph{columns} are the spherical basis vectors expressed in
Cartesian components is
\begin{equation}
  R(\theta,\phi)
  \;=\;
  \begin{pmatrix}
    \sin\theta\cos\phi & \cos\theta\cos\phi & -\sin\phi \\
    \sin\theta\sin\phi & \cos\theta\sin\phi &  \cos\phi \\
    \cos\theta         & -\sin\theta        &  0
  \end{pmatrix},
  \label{eq:R-matrix}
\end{equation}
so that a vector $\mathbf{v}$ with spherical components
$(v_r, v_\theta, v_\phi)^T$ has Cartesian components
\begin{equation}
  \begin{pmatrix}v_x\\v_y\\v_z\end{pmatrix}
  = R(\theta,\phi)
  \begin{pmatrix}v_r\\v_\theta\\v_\phi\end{pmatrix}.
  \label{eq:sph-to-cart-vec}
\end{equation}
Since $R$ is orthogonal ($R^T = R^{-1}$, verified by direct multiplication
$R^T R = I_3$), the inverse is
\begin{equation}
  \begin{pmatrix}v_r\\v_\theta\\v_\phi\end{pmatrix}
  = R^T(\theta,\phi)
  \begin{pmatrix}v_x\\v_y\\v_z\end{pmatrix}.
  \label{eq:cart-to-sph-vec}
\end{equation}

\begin{remark}[Rotating frame]
  The spherical frame $(\hat{e}_r,\hat{e}_\theta,\hat{e}_\phi)$ rotates
  with the satellite's position: as $\theta$ and $\phi$ change along the
  orbit, $R$ changes.  When converting the spherical-frame acceleration
  components back to inertial Cartesian, apply~\eqref{eq:sph-to-cart-vec}
  evaluated at the \emph{current} $(r,\theta,\phi)$.
\end{remark}

% ---------------------------------------------------------------------------
\subsection{Gravitational Potential: Multipole Expansion from First Principles}
\label{ssec:multipole}
% ---------------------------------------------------------------------------

The Newtonian gravitational potential at exterior point $\mathbf{r}$ due to a
body with mass density $\rho(\mathbf{r}')$ is
\begin{equation}
  V(\mathbf{r})
  = -G\int_{\mathcal{B}} \frac{\rho(\mathbf{r}')}{|\mathbf{r}-\mathbf{r}'|}\,dV'.
  \label{eq:V-integral}
\end{equation}
(Convention: $V < 0$ for bound orbits.  Some references use $U = -V > 0$.)

\paragraph{Legendre generating function.}
For $r = |\mathbf{r}| > r' = |\mathbf{r}'|$ (exterior expansion), the identity
\begin{equation}
  \frac{1}{|\mathbf{r}-\mathbf{r}'|}
  = \frac{1}{r}\sum_{n=0}^\infty
    \left(\frac{r'}{r}\right)^n P_n(\cos\gamma),
  \label{eq:legendre-gen}
\end{equation}
where $\gamma$ is the angle between $\mathbf{r}$ and $\mathbf{r}'$, follows from
the generating function of Legendre polynomials
$\sum_{n=0}^\infty t^n P_n(u) = (1-2tu+t^2)^{-1/2}$ with $t=r'/r$,
$u=\cos\gamma$.  The series converges absolutely for $r > r'$.

Substituting~\eqref{eq:legendre-gen} into~\eqref{eq:V-integral}:
\begin{equation}
  V(\mathbf{r})
  = -\frac{G}{r}\sum_{n=0}^\infty \frac{1}{r^n}
    \int_{\mathcal{B}} \rho(\mathbf{r}')\,(r')^n\,P_n(\cos\gamma)\,dV'.
  \label{eq:V-series-gamma}
\end{equation}

\paragraph{Spherical harmonic addition theorem.}
The addition theorem for $P_n$ decouples the angles of $\mathbf{r}$ and
$\mathbf{r}'$:
\begin{equation}
  P_n(\cos\gamma)
  = \frac{4\pi}{2n+1}\sum_{m=-n}^{n}
    Y_{nm}^*(\theta',\phi')\,Y_{nm}(\theta,\phi),
  \label{eq:addition-theorem}
\end{equation}
where $Y_{nm}$ are the complex spherical harmonics.  Substituting into
\eqref{eq:V-series-gamma} and integrating over $\mathbf{r}'$, the
\emph{multipole moments} of the body appear naturally:
\begin{equation}
  V(\mathbf{r})
  = -G\sum_{n=0}^\infty\sum_{m=-n}^{n}
    \frac{4\pi}{(2n+1)\,r^{n+1}}
    \!\left[\int_{\mathcal{B}}\rho(\mathbf{r}')\,(r')^n\,Y_{nm}^*(\theta',\phi')\,dV'\right]
    Y_{nm}(\theta,\phi).
  \label{eq:V-multipole-complex}
\end{equation}

\paragraph{Real harmonic form.}
Converting to real tesseral harmonics via
$Y_{nm}^c = \sqrt{2}\,\mathrm{Re}[Y_{nm}]$ and $Y_{nm}^s = \sqrt{2}\,\mathrm{Im}[Y_{nm}]$,
and introducing the reference equatorial radius $R_e$ and total mass $M$ to
non-dimensionalise the coefficients, one obtains the standard
\emph{geopotential expansion}:
\begin{equation}
  \boxed{
  V(r,\theta,\phi)
  = \frac{\mu}{r}\sum_{n=0}^\infty\sum_{m=0}^{n}
    \left(\frac{R_e}{r}\right)^{\!n}
    P_{nm}(\cos\theta)
    \Bigl[C_{nm}\cos(m\phi) + S_{nm}\sin(m\phi)\Bigr]
  }
  \label{eq:geopotential}
\end{equation}
where $\mu = GM$ is the standard gravitational parameter, $R_e$ is the mean
equatorial radius, $P_{nm}$ are the \emph{associated Legendre functions}
(see Section~\ref{ssec:legendre}), and $C_{nm}$, $S_{nm}$ are the dimensionless
Stokes coefficients (gravity field harmonics).

\begin{remark}[Sign convention]
  Some references (Kaula~\cite{Kaula1966}; older geodesy) define $V$ with a
  positive sign (force potential $U = -V$).  The convention
  in~\eqref{eq:geopotential} gives $V < 0$ and acceleration
  $\mathbf{a} = -\nabla V$.
\end{remark}

% ---------------------------------------------------------------------------
\subsection{Legendre Polynomials and Associated Legendre Functions}
\label{ssec:legendre}
% ---------------------------------------------------------------------------

The \emph{Legendre polynomials} $P_n(u)$ are the $m=0$ solutions of Legendre's
differential equation $(1-u^2)P_n'' - 2uP_n' + n(n+1)P_n = 0$.  The first few
are:
\begin{align}
  P_0(u) &= 1,\label{eq:P0}\\
  P_1(u) &= u,\label{eq:P1}\\
  P_2(u) &= \tfrac{1}{2}(3u^2-1),\label{eq:P2}\\
  P_3(u) &= \tfrac{1}{2}(5u^3-3u),\label{eq:P3}\\
  P_4(u) &= \tfrac{1}{8}(35u^4-30u^2+3).\label{eq:P4}
\end{align}
They satisfy the orthogonality relation
$\int_{-1}^1 P_n(u)\,P_{n'}(u)\,du = \tfrac{2}{2n+1}\delta_{nn'}$.

The \emph{associated Legendre functions} (un-normalized, matching
\eqref{eq:geopotential}) are defined by
\begin{equation}
  P_{nm}(u) = (1-u^2)^{m/2}\,\frac{d^m}{du^m}P_n(u),
  \quad n \geq 0,\; 0 \leq m \leq n.
  \label{eq:Pnm-def}
\end{equation}
Explicitly for the terms needed through $n=3$:
\begin{align}
  P_{20}(u) &= P_2(u) = \tfrac{1}{2}(3u^2-1),\label{eq:P20}\\
  P_{21}(u) &= 3u\sqrt{1-u^2},\label{eq:P21}\\
  P_{22}(u) &= 3(1-u^2),\label{eq:P22}\\
  P_{30}(u) &= P_3(u) = \tfrac{1}{2}(5u^3-3u),\label{eq:P30}\\
  P_{31}(u) &= \tfrac{3}{2}(5u^2-1)\sqrt{1-u^2},\label{eq:P31}\\
  P_{32}(u) &= 15u(1-u^2),\label{eq:P32}\\
  P_{33}(u) &= 15(1-u^2)^{3/2}.\label{eq:P33}
\end{align}

\paragraph{Harmonic taxonomy.}
The $(n,m)$ harmonics are classified as:
\begin{itemize}
  \item \textbf{Zonal} ($m=0$): depend only on latitude; axisymmetric rings.
        Coefficient $J_n \equiv -C_{n0}$.  Earth's oblateness is dominated by
        $J_2 \approx 1.08\times10^{-3}$.
  \item \textbf{Tesseral} ($0 < m < n$): depend on both latitude and longitude;
        ``checkerboard'' patterns on the sphere.
  \item \textbf{Sectoral} ($m=n$): depend primarily on longitude;
        ``orange-peel'' slices.
\end{itemize}

% ---------------------------------------------------------------------------
\subsection{Explicit Expansion: Monopole through $J_3$}
\label{ssec:explicit}
% ---------------------------------------------------------------------------

We expand \eqref{eq:geopotential} term by term, substituting $u = \cos\theta$
and the Cartesian relations \eqref{eq:sph2cart-x}--\eqref{eq:sph2cart-z}.

\subsubsection{$n=0$: Monopole (Point-Mass)}

\begin{equation}
  V_0 = \frac{\mu}{r},
  \label{eq:V0}
\end{equation}
with $C_{00} = 1$ by normalization (total-mass convention).
This is the two-body Keplerian potential.

\subsubsection{$n=1$: Dipole (Vanishes at Centre of Mass)}

The $n=1$ terms, using $P_{10}(u)=u$ and $P_{11}(u)=\sqrt{1-u^2}$, are:
\begin{equation}
  V_1 = \frac{\mu R_e}{r^3}
        \Bigl[C_{10}\,z + C_{11}\,x + S_{11}\,y\Bigr].
  \label{eq:V1-cart}
\end{equation}
The coefficients $C_{10}$, $C_{11}$, $S_{11}$ are proportional to the
components of the body's centre of mass expressed in the body-fixed frame.
\emph{When the origin is chosen at the body's centre of mass,
$C_{10}=C_{11}=S_{11}=0$}, and $V_1 \equiv 0$.  This is the standard
assumption for geocentric orbit mechanics.

\subsubsection{$n=2$, $m=0$: The $J_2$ Zonal Term}

\begin{equation}
  V_{20} = \frac{\mu}{r}\left(\frac{R_e}{r}\right)^{\!2}
            P_2(\cos\theta)\,C_{20}.
  \label{eq:V20}
\end{equation}
Using \eqref{eq:P2} with $u=\cos\theta=z/r$:
\begin{equation}
  P_2(\cos\theta)
  = \tfrac{1}{2}(3\cos^2\theta - 1)
  = \tfrac{1}{2}\!\left(\frac{3z^2}{r^2} - 1\right)
  = \frac{2z^2 - x^2 - y^2}{2r^2}.
  \label{eq:P2-cart}
\end{equation}
Defining $J_2 = -C_{20} > 0$ (Earth is oblate, so $C_{20} < 0$):
\begin{equation}
  \boxed{
  V_{J_2}
  = -\frac{\mu J_2 R_e^2}{2r^3}\left(3\cos^2\theta - 1\right)
  = -\frac{\mu J_2 R_e^2}{2r^5}\left(2z^2 - x^2 - y^2\right).
  }
  \label{eq:V-J2}
\end{equation}

The $n=2$ tesseral ($m=1$) and sectoral ($m=2$) harmonics are:
\begin{align}
  V_{21} &= \frac{3\mu R_e^2}{r^5}
             \bigl(C_{21}\,xz + S_{21}\,yz\bigr),
  \label{eq:V21}\\
  V_{22} &= \frac{3\mu R_e^2}{r^5}
             \bigl[C_{22}(x^2-y^2) + 2S_{22}\,xy\bigr].
  \label{eq:V22}
\end{align}
For Earth, $|C_{21}|,|S_{21}| < 10^{-9}$ (negligible);
$C_{22} \approx 1.574\times10^{-6}$ encodes the equatorial ellipticity.

\subsubsection{$n=3$, $m=0$: The $J_3$ Zonal Term}

\begin{equation}
  V_{30} = \frac{\mu}{r}\left(\frac{R_e}{r}\right)^{\!3}
            P_3(\cos\theta)\,C_{30},
  \quad J_3 = -C_{30}.
  \label{eq:V30}
\end{equation}
Using \eqref{eq:P3} with $u=\cos\theta=z/r$:
\begin{equation}
  P_3(\cos\theta)
  = \tfrac{1}{2}(5\cos^3\theta - 3\cos\theta)
  = \frac{z(5z^2 - 3r^2)}{2r^3}.
  \label{eq:P3-cart}
\end{equation}
Therefore:
\begin{equation}
  \boxed{
  V_{J_3}
  = -\frac{\mu J_3 R_e^3}{2r^7}\,z\bigl(5z^2 - 3r^2\bigr).
  }
  \label{eq:V-J3}
\end{equation}
For Earth, $J_3 \approx -2.532\times10^{-6}$ (negative: slight
southern-hemisphere mass excess, the ``pear-shape'' perturbation).

\paragraph{Numerical constants (WGS84 / EGM96).}
\begin{align*}
  \mu    &= 3.986\,004\,418\times10^{14}\ \mathrm{m^3\,s^{-2}},\\
  R_e    &= 6\,378\,137\ \mathrm{m},\\
  J_2    &= 1.082\,63\times10^{-3},\\
  J_3    &\approx -2.532\times10^{-6}.
\end{align*}
The full zonal potential through $J_3$ in compact spherical form is:
\begin{equation}
  V \approx \frac{\mu}{r}\left[
    1
    - J_2\left(\frac{R_e}{r}\right)^{\!2}\!\frac{3\cos^2\theta-1}{2}
    - J_3\left(\frac{R_e}{r}\right)^{\!3}\!\frac{5\cos^3\theta-3\cos\theta}{2}
  \right].
  \label{eq:V-J2J3}
\end{equation}

% ---------------------------------------------------------------------------
\subsection{Gradient and Acceleration}
\label{ssec:gradient}
% ---------------------------------------------------------------------------

The gravitational acceleration is $\mathbf{a} = -\nabla V$.  The gradient in
spherical coordinates is
\begin{equation}
  \nabla V
  = \frac{\partial V}{\partial r}\hat{e}_r
  + \frac{1}{r}\frac{\partial V}{\partial \theta}\hat{e}_\theta
  + \frac{1}{r\sin\theta}\frac{\partial V}{\partial \phi}\hat{e}_\phi.
  \label{eq:grad-sph}
\end{equation}

\subsubsection{Two-Body (Keplerian) Acceleration}

\begin{equation}
  \mathbf{a}_0 = -\frac{\mu}{r^2}\hat{e}_r
  = -\frac{\mu}{r^3}\mathbf{r},
  \label{eq:a0}
\end{equation}
so in Cartesian: $\mathbf{a}_0 = -(\mu/r^3)(x,y,z)^T$.

\subsubsection{$J_2$ Acceleration in Spherical Coordinates}

Differentiating \eqref{eq:V-J2} with respect to $r$ and $\theta$:
\begin{align}
  \frac{\partial V_{J_2}}{\partial r}
  &= \frac{3\mu J_2 R_e^2}{2r^4}(3\cos^2\theta-1),
  \label{eq:dVJ2dr}\\
  \frac{1}{r}\frac{\partial V_{J_2}}{\partial\theta}
  &= \frac{3\mu J_2 R_e^2}{r^4}\cos\theta\sin\theta.
  \label{eq:dVJ2dtheta}
\end{align}
Since $V_{J_2}$ is independent of $\phi$, the $\hat{e}_\phi$ component is zero.
The spherical-frame $J_2$ acceleration is:
\begin{align}
  a_{J_2,r}      &= -\frac{3\mu J_2 R_e^2}{2r^4}(3\cos^2\theta-1),
  \label{eq:aJ2r}\\
  a_{J_2,\theta} &= -\frac{3\mu J_2 R_e^2}{r^4}\cos\theta\sin\theta,
  \label{eq:aJ2theta}\\
  a_{J_2,\phi}   &= 0.
  \label{eq:aJ2phi}
\end{align}

\subsubsection{$J_2$ Acceleration in Cartesian Coordinates}

Apply the rotation~\eqref{eq:sph-to-cart-vec}:
\begin{equation}
  \begin{pmatrix}a_x\\a_y\\a_z\end{pmatrix}_{J_2}
  = R(\theta,\phi)
  \begin{pmatrix}a_{J_2,r}\\a_{J_2,\theta}\\0\end{pmatrix}.
  \label{eq:aJ2-rotate}
\end{equation}
Expanding with $\cos\theta = z/r$,
$\sin\theta\cos\phi = x/r$, $\sin\theta\sin\phi = y/r$:
\begin{align}
  a_x^{J_2} &= a_{J_2,r}\sin\theta\cos\phi + a_{J_2,\theta}\cos\theta\cos\phi
             = \cos\phi\bigl(a_{J_2,r}\sin\theta + a_{J_2,\theta}\cos\theta\bigr).
  \label{eq:aJ2x-expand}
\end{align}
Substituting \eqref{eq:aJ2r}--\eqref{eq:aJ2theta} and simplifying, the result
is the compact form used in orbital mechanics codes:
\begin{equation}
  \boxed{
  \mathbf{a}_{J_2}
  = -\frac{3\mu J_2 R_e^2}{2r^5}
    \begin{pmatrix}
      x\!\left(1 - 5z^2/r^2\right)\\[2pt]
      y\!\left(1 - 5z^2/r^2\right)\\[2pt]
      z\!\left(3 - 5z^2/r^2\right)
    \end{pmatrix}.
  }
  \label{eq:aJ2-compact}
\end{equation}

\begin{remark}[Equatorial check]
  Setting $z=0$ (equatorial orbit): $\mathbf{a}_{J_2} = -(3\mu J_2 R_e^2/2r^5)
  (x,y,0)^T$.  This is a radially outward perturbation in the equatorial
  plane, consistent with the oblate equatorial bulge attracting a test mass
  less strongly than a sphere of the same mass.
\end{remark}

\begin{remark}[Polar check]
  Setting $x=y=0$, $z=r$ (polar pass): $\mathbf{a}_{J_2} = -(3\mu J_2 R_e^2/2r^5)
  \cdot z(3-5) \hat{z} = +(\mu J_2 R_e^2/r^5)z\,\hat{z}$.  The $J_2$
  acceleration is directed \emph{away} from Earth at the poles, as expected
  (less equatorial mass overhead means slightly less gravitational pull).
\end{remark}

\subsubsection{$J_3$ Acceleration in Cartesian Coordinates}

Differentiating \eqref{eq:V-J3} and applying the same rotation procedure:
\begin{equation}
  \boxed{
  \mathbf{a}_{J_3}
  = -\frac{5\mu J_3 R_e^3}{2r^7}
    \begin{pmatrix}
      xz\!\left(3 - 7z^2/r^2\right)\\[2pt]
      yz\!\left(3 - 7z^2/r^2\right)\\[2pt]
      \frac{3r^2}{5} - 6z^2 + \frac{7z^4}{r^2}
    \end{pmatrix}.
  }
  \label{eq:aJ3-compact}
\end{equation}
Note that $a_x^{J_3}$ and $a_y^{J_3}$ are both proportional to $z$: the
$J_3$ perturbation has no in-plane component at the equator ($z=0$), only
north--south asymmetry.

% ---------------------------------------------------------------------------
\subsection{Summary: Total Acceleration Through $J_3$}
\label{ssec:gravity-summary}
% ---------------------------------------------------------------------------

The total gravitational acceleration in inertial Cartesian (ECI) coordinates,
including the two-body term and zonal harmonics through $J_3$, is:
\begin{equation}
  \mathbf{a}(\mathbf{r})
  = -\frac{\mu}{r^3}\mathbf{r}
  + \mathbf{a}_{J_2}
  + \mathbf{a}_{J_3}
  + \mathcal{O}(J_4),
  \label{eq:a-total}
\end{equation}
with $\mathbf{a}_{J_2}$ from \eqref{eq:aJ2-compact} and $\mathbf{a}_{J_3}$
from \eqref{eq:aJ3-compact}.  For a low-Earth orbit satellite,
$|\mathbf{a}_{J_2}|/|\mathbf{a}_0| \approx J_2(R_e/r)^2 \approx 10^{-3}$
while $|\mathbf{a}_{J_3}|/|\mathbf{a}_0| \approx 10^{-6}$.

This acceleration is the right-hand side of the second-order ODE
\eqref{eq:stormer-ode} propagated by the St\"{o}rmer/Numerov integrators
described in the preceding sections.

\end{document}
